\section*{الفصل \ref{ch:b2:schrodinger-wave-functions} --- معادلة شرودنغر والتوابع الموجية}

\begin{solution}{exo:b2:schrodinger-wave-functions:1}
$h/\sqrt{2mE}$: \qty{0.123}{nm}، \qty{12.3}{pm}، \qty{3.9}{pm} (وعند \qty{100}{keV} تكون
الطاقة الحركية خُمس $mc^2$؛ وكمية الحركة النسبية تعطي
\qty{3.7}{pm})؛ والنترون \qty{0.18}{nm}؛ و C$_{60}$ \qty{2.8}{pm}؛ وكرة التنس
$\sim 10^{-34}\,\unit{m}$. وكلها إلا الكرة أظهرت تداخلًا.
\end{solution}

\begin{solution}{exo:b2:schrodinger-wave-functions:2}
$A = (\pi a^2)^{-1/4}$؛ و $\langle x\rangle = 0$، و $\Delta x = a/\sqrt2$؛ و $0.84$؛ و $0.9995$.
\end{solution}

\begin{solution}{exo:b2:schrodinger-wave-functions:3}
التعويض يعطي $E\varphi = -(\hbar^2/2m)\varphi'' + V\varphi$. و $E/h =
\qty{2.4e14}{Hz}$؛ و \qty{4.8e14}{Hz}، و \qty{620}{nm}: أي الضوء الأحمر الذي يصدره
الانتقال.
\end{solution}

\begin{solution}{exo:b2:schrodinger-wave-functions:4}
الإلكترون، \qty{1}{nm}: $\tau = \qty{1.7e-14}{s}$، و \qty{60}{\micro m} بعد \qty{1}{ns}،
و \qty{60}{km} بعد \qty{1}{s}؛ ومع \qty{1}{mm}: $\tau = \qty{17}{s}$، بلا تغير بعد \qty{1}{ns}،
و $+0.2\%$ بعد \qty{1}{s}؛ والبروتون، \qty{1}{nm}: $\tau = \qty{3.2e-11}{s}$، و \qty{30}{nm}،
و \qty{30}{m}؛ والحبيبة: $\tau = 600$ سنة. وللأجسام المرئية يكون $\tau$
فلكيًّا --- ومحيطها يعيد توضيعها باستمرار.
\end{solution}

\begin{solution}{exo:b2:schrodinger-wave-functions:5}
(أ) \cref{prop:b2:schrodinger-wave-functions:current}. (ب) و $|A|^2\hbar k/m$؛
و $0$؛ و $(|A|^2 - |B|^2)\hbar k/m$. (ج) والتدفقان الوارد والمنعكس
يتطارحان بلا حدّ تداخل: فيمكن تعريف معاملي النفاذ
والانعكاس من التدفقات. (د) و $\psi^*\psi'$ حقيقي، فجزؤه
التخيلي معدوم.
\end{solution}

\begin{solution}{exo:b2:schrodinger-wave-functions:6}
(أ) $A = (2\pi\sigma^2)^{-1/4}$، و $\langle x\rangle = 0$، و $\Delta x = \sigma$. (ب) و $\hbar k_0$.
(ج) و $\sigma \times \hbar/2\sigma = \hbar/2$. (د) و $\dd\langle x\rangle/\dd t = \int x\,
\partial_t\rho = -\int x\,\partial_xj = \int j = (\hbar/m)\operatorname{Im}\int\psi^*
\psi' = \langle p\rangle/m$.
\end{solution}

\begin{solution}{exo:b2:schrodinger-wave-functions:7}
(أ) \qty{5.9e6}{m/s}، و \qty{0.123}{nm}، و \qty{51}{ns}. (ب) و $v_\varphi = \qty{3.0e6}{m/s}$،
و $v_{\text{g}} = \qty{5.9e6}{m/s}$: فالأخيرة. (ج) و $\tau = \qty{1.7}{\micro s} \gg
\qty{51}{ns}$: أي بلا تغير. (د) و $\lambda/a = \qty{4e-7}{rad}$، و \qty{0.1}{\micro m}:
فالأشعة صالحة.
\end{solution}

\begin{solution}{exo:b2:schrodinger-wave-functions:8}
(أ) $|\psi|^2$ وكل قيمة متوقعة لا تتغير. (ب) و $\tfrac12(|\psi_1|^2
+ |\psi_2|^2) + \operatorname{Re}(\eu^{\iu\alpha}\psi_1^*\psi_2)$. (ج) وهما الموجتان الآتيتان من
الشقّين؛ وطورهما النسبي $k(r_2 - r_1)$. (د) وينتهي الكاشف
إلى حالتين مختلفتين من أجل المسارين؛ فلا تتداخل المركّبتان
بعد ذلك، ويكون النقش مجموع نقشي شقّ
واحد.
\end{solution}

\begin{solution}{exo:b2:schrodinger-wave-functions:9}
(أ) $\langle V'\rangle = m\omega^2\langle x\rangle$: أي بالضبط. (ب) و $\langle V'\rangle
= mg$: أي بالضبط. (ج) وحين تكون الرزمة ضيقة في وجه السلّم الذي
تتغير عليه $V''$. (د) وطول موجتها عند \qty{0.1}{m/s} هو $10^{-29}\,\unit{m}$؛ وكانت
الرزمة يجب أن تكون بعرض الوعاء لتهمّ. وطول موجة الذرة
الحراري، \qty{0.1}{nm}، هو حجمها نفسه.
\end{solution}

\begin{solution}{exo:b2:schrodinger-wave-functions:10}
(أ) $\tfrac12(\varphi_1^2 + \varphi_2^2) + \varphi_1\varphi_2\cos\omega_{21}t$. (ب) و $L/2 +
\cos\omega_{21}t\int x\varphi_1\varphi_2 = L/2 - (16L/9\pi^2)\cos\omega_{21}t$. (ج)
و $3h^2/8mL^2 = \qty{1.13}{eV}$، و \qty{2.7e14}{Hz}، و \qty{1.1}{\micro m}. (د) وهي تشعّ
عند $\nu_{21}$: أي الخط؛ وللحالة المستقرة الصرفة لا \omterm{thm:b2:dipole-radiation:field}{ثنائي قطب
متذبذب} فلا تشعّ.
\end{solution}

\begin{solution}{exo:b2:schrodinger-wave-functions:11}
(أ) $\psi \propto \int\eu^{-\sigma_0^2k^2}\eu^{\iu(kx - \hbar k^2t/2m)}\dd k$. (ب) و $\alpha =
\sigma_0^2 + \iu\hbar t/2m$، و $\beta = \iu x$: $|\psi|^2 \propto \exp[-x^2/2\sigma_0^2(1 +
(\hbar t/2m\sigma_0^2)^2)]$. (ج) و $\Delta v = \hbar/2m\sigma_0$، و $\Delta v\,\tau =
\sigma_0$. (د) و $\omega = ck$: فلا تبدد، وكل المركّبات عند $c$.
\end{solution}

\begin{solution}{exo:b2:schrodinger-wave-functions:12}
(أ) $\Gamma = \hbar/\tau$: \qty{6.6e-8}{eV}؛ و \qty{6.6e-8}{eV}؛ و \qty{66}{MeV}. (ب)
و $1/2\pi\tau = \qty{16}{MHz}$، أي أقلّ بمئة مرة من دوبلر. (ج) وفي
الذرات الباردة والشوارد المحبوسة يزول عرض دوبلر ويُبلَغ العرض
الطبيعي --- وهو أساس معايير التواتر. (د) و $c\tau =
\qty{3}{m}$، أي طول القطار.
\end{solution}

\begin{solution}{pb:b2:schrodinger-wave-functions:1}
\textbf{1.} $v = \sqrt{2eU/m} = \qty{8.4e7}{m/s}$؛ و $p = \qty{7.6e-23}{kg.m/s}$؛ و $\lambda =
\qty{8.7}{pm}$.

\textbf{2.} أخطاء بضعة في المئة: أي مقبولة في التقديرات (و $\lambda =
\qty{8.6}{pm}$ بالضبط).

\textbf{3.} \qty{4.8}{ns}.

\textbf{4.} $\lambda/a = \qty{3e-8}{rad}$: أي \qty{12}{nm} --- أصغر $10^4$ مرة من
بكسل.

\textbf{5.} $\Delta\theta = \Delta p_y/p \sim \hbar/2ap = \lambda/4\pi a$: أي الرتبة نفسها.

\textbf{6.} $\tau = \qty{1.6}{ms} \gg \qty{4.8}{ns}$: أي بلا تغير.

\textbf{7.} $\dd\langle p_y\rangle/\dd t = eE$ بالضبط من أجل حقل منتظم:
فتتبع $\langle y\rangle$ القطع المكافئ الكلاسيكي؛ والأنبوب أداة
كلاسيكية.

\textbf{8.} $v/2 = \qty{4.2e7}{m/s}$؛ ولا تُرصَد إلا \omterm{prop:b2:dispersion-wave-packets:group}{سرعة المجموعة}؛
وطور الموجة المستوية لا يحمل إشارة (ويتعلق بالصفر
الاعتباطي للطاقة).

\textbf{9.} $\Delta p = \hbar/2\sigma_0$، و $\Delta v = \hbar/2m\sigma_0$.

\textbf{10.} \qty{1.7e-16}{s}، \qty{1.7e-14}{s}، \qty{1.7e-8}{s}؛ والبروتون \qty{3.2e-11}{s}؛
والفيروس \qty{20}{s}.

\textbf{11.} لأنه في \omterm{thm:b2:schrodinger-wave-functions:stationary}{حالة مستقرة}: فالكمون يمسك
الرزمة؛ و $|\varphi|^2$ لا يتحرك.

\textbf{12.} $v = \qty{5.9e5}{m/s}$؛ والتضاعف عند $t = \sqrt3\tau = \qty{2.9e-14}{s}$:
أي \qty{17}{nm}، ونحو أربعة عشر طول موجة.

\textbf{13.} توابعها الموجية تمتد على البلورة كلها: أي كمية
حركة معيّنة، ولا موضع معيّن.

\textbf{14.} التمييز $\sim\lambda/\alpha = \qty{4}{pm}/0.01 = \qty{0.4}{nm}$: فزيغانات
العدسات المغناطيسية تحدّ الفتحة، لا طول الموجة.

\textbf{15.} قارن $\lambda$ بالأحجام الداخلة في اللعب (كالفتحات، وسلّم
$V$)، وقارن زمن الانتشار بزمن الرصد.

\textbf{16.} \qty{0.38}{}، \qty{1.50}{}، \qty{3.38}{eV}؛ و $\nu_{21} = \qty{2.7e14}{Hz}$
(\qty{1.1}{\micro m})، و $\nu_{31} = \qty{7.3e14}{Hz}$ (\qty{410}{nm}).

\textbf{17.} \omterm{thm:b2:gratings:nwaves}{للشعاع الطوري} طويلة واحدة؛ وعند $x = L/2$.

\textbf{18.} عند $t = 0$ تتكدس الكثافة في النصف الأيسر، وعند
$T_{21}/2$ في الأيمن.

\textbf{19.} $\langle x\rangle = L/2 - (16L/9\pi^2)\cos\omega_{21}t$: أي سعة
$0.18L = \qty{0.18}{nm}$.

\textbf{20.} عند $\nu_{21}$: أي الخط الطيفي؛ وللحالة المستقرة لا \omterm{thm:b2:dipole-radiation:field}{ثنائي
قطب متذبذب} فهي مستقرة؛ والتراكب يشعّ حتى ينزل
إلى $\varphi_1$.

\textbf{21.} ليست معيّنة: $\langle E\rangle = (E_1 + E_2)/2 = \qty{0.94}{eV}$؛ و
القياس يعطي $E_1$ أو $E_2$، وكلٌّ باحتمال النصف.

\textbf{22.} العشوائي: أين تقع كل إلكترون؛ وغير العشوائي:
النقش $|\psi|^2$ الذي تبنيه المواقع.

\textbf{23.} $j = |A|^2\hbar k/m$: فمع $|A|^2 = n$ إلكترونًا في وحدة الطول،
يكون $nv$ إلكترونًا في الثانية.

\textbf{24.} من المرتبة الأولى حتى يكون $\psi(x,0)$ وصفًا كاملًا
يثبّت المستقبل؛ وعقدية حتى يكون للطاقة المعيّنة التعلق
الزمني الوحيد $\eu^{-\iu Et/\hbar}$ وحتى يخرج $E = p^2/2m$.

\textbf{25.} $\psi$ هي السعة التي مربعها الاحتمال؛ وهي
تحقق معادلة شرودنغر الخطية؛ \omterm{thm:b2:schrodinger-wave-functions:stationary}{والحالات المستقرة} تحمل
شعاعًا طوريًّا وكثافات ثابتة؛ والرزمة الحرة تتحرك بالسرعة $p/m$ وتنتشر
في $2m\sigma_0^2/\hbar$؛ وإهرنفست يعيد نيوتن من أجل الكبير.
\end{solution}
